% SmartMess Product Requirements Document
% Formatting modeled after the supplied reference PDF:
% "Grievance-Resolution System for a Tier-2/3 City"
%
% The SmartMess content below is based on the supplied SmartMess_PRD.docx.
% Replace the metadata placeholders below if course/instructor/author details
% are required for your submission.

\documentclass[12pt,a4paper]{article}

\usepackage[a4paper,margin=0.78in,top=0.72in,bottom=0.72in]{geometry}
\usepackage[T1]{fontenc}
\usepackage[utf8]{inputenc}
\usepackage{lmodern}
\usepackage{microtype}
\usepackage{enumitem}
\usepackage{array}
\usepackage{tabularx}
\usepackage{booktabs}
\usepackage{longtable}
\usepackage{ragged2e}
\usepackage{amsmath}
\usepackage{fancyhdr}
\usepackage[hidelinks]{hyperref}

% ---------- General formatting ----------
\setlength{\parindent}{0pt}
\setlength{\parskip}{0.55em}
\setlength{\emergencystretch}{3em}

\setlist[itemize]{
    leftmargin=1.35em,
    itemsep=0.35em,
    topsep=0.25em,
    parsep=0pt
}

\setlist[enumerate]{
    leftmargin=1.55em,
    itemsep=0.35em,
    topsep=0.25em,
    parsep=0pt
}

% Match the compact numbered-section appearance of the reference.
\usepackage{titlesec}

\titleformat{\section}
  {\normalfont\large\bfseries}
  {\thesection}
  {0.65em}
  {}

\titleformat{\subsection}
  {\normalfont\normalsize\bfseries}
  {\thesubsection}
  {0.65em}
  {}

\titlespacing*{\section}
  {0pt}{1.25em}{0.45em}

\titlespacing*{\subsection}
  {0pt}{0.85em}{0.35em}

% ---------- Header / footer ----------
\pagestyle{fancy}
\fancyhf{}
\renewcommand{\headrulewidth}{0pt}
\renewcommand{\footrulewidth}{0pt}
\fancyfoot[C]{\thepage}

% ---------- Helper commands ----------
\newcommand{\annotation}[1]{\hfill\textit{---#1}}

\newcommand{\priority}[1]{\textbf{[#1]}}

\newcolumntype{Y}{>{\RaggedRight\arraybackslash}X}
\newcolumntype{L}[1]{>{\RaggedRight\arraybackslash}p{#1}}

% Prevent widows/orphans where possible.
\widowpenalty=10000
\clubpenalty=10000

\begin{document}

% ============================================================
% TITLE AREA
% ============================================================

\begin{center}
    {\LARGE\bfseries SmartMess\par}
    \vspace{0.35em}
    {\large Product Requirements Document\par}
    \vspace{0.15em}
    {\normalsize Consumption-aware mess billing \& operations platform\par}

    % Optional academic metadata. Fill these only if required.
    % \vspace{1.2em}
    % Submitted to: [Instructor Name]\\
    % [Institution Name]\\
    % [Author Name(s)]\\
    % [Date]
\end{center}

\vspace{0.8em}

% ============================================================
% EXECUTIVE SUMMARY
% ============================================================

\section*{Executive Summary}

Hostel mess programs typically charge a fixed fee regardless of how many meals a
student actually eats, while mess operators plan food production from estimated
rather than confirmed demand. This disconnect causes students to overpay for
unused meals and causes operators to over- or under-prepare food.

SmartMess is a consumption-aware mess management platform that links what a
student actually eats to what they are billed and to how much food the mess
prepares. Students record meal participation and see transparent, usage-based
charges; administrators get a single system for billing, menus, feedback, and
demand planning.

The MVP is scoped in three phases --- consumption \& billing transparency,
administration \& feedback, and food-efficiency analytics --- so that each phase
validates a distinct hypothesis before the next is built.

% ============================================================
% 1. PROBLEM & OPPORTUNITY
% ============================================================

\section{Problem \& Opportunity}

\subsection*{Problem}

Traditional hostel mess systems charge a fixed fee irrespective of actual
consumption. A student may be away from campus, eating outside, home for the
weekend, skipping meals, or absent for classes or events --- and still pay the
same amount.

Mess operators face a mirrored problem: food is prepared against estimated or
historical demand rather than confirmed meal participation, so low turnout
produces excess food and unexpected turnout produces shortages.

\subsection*{Core insight}

Both problems trace back to the same root cause --- the mess has no reliable,
real-time signal of who is actually eating. Fixing consumption visibility fixes
billing fairness and food planning at the same time.

\subsection*{The Opportunity}

SmartMess replaces two independent, low-visibility processes with one shared
data model of meal consumption:

\begin{center}
\begin{tabularx}{0.92\textwidth}{L{0.42\textwidth} Y}
\toprule
\textbf{From} & \textbf{To}\\
\midrule
Fixed fee, regardless of consumption &
Meal usage $\rightarrow$ recorded consumption $\rightarrow$ fairer billing\\
Estimated demand &
Meal participation $\rightarrow$ better demand estimation $\rightarrow$ less food waste\\
\bottomrule
\end{tabularx}
\end{center}

% ============================================================
% 2. TARGET USERS & USE CASES
% ============================================================

\section{Target Users \& Use Cases}

\subsection*{Primary user --- Student}

Students are the main beneficiaries of the platform.

\textbf{Key use cases}
\begin{itemize}
    \item View available meals and the daily/weekly menu.
    \item Indicate whether they will consume a meal, or opt out when absent.
    \item Track meals consumed and accumulated mess usage/cost.
    \item Pay based on actual eligible consumption, where the institution supports
          usage-based billing.
    \item Provide meal feedback and report food or service problems.
\end{itemize}

\textbf{Core student value}

Pay more fairly, understand your mess usage, and have a voice in the quality of
the food you receive.

\subsection*{Secondary user --- Mess Administrator}

Administrators manage the operational and financial side of the mess.

\textbf{Key use cases}
\begin{itemize}
    \item Manage student accounts and meal rates/billing rules.
    \item Monitor meal participation and track actual consumption.
    \item Manage menus, view feedback, and resolve complaints.
    \item Monitor expected food demand and food wastage.
    \item Generate operational and billing reports.
\end{itemize}

\subsection*{Stakeholder --- Hostel / College Management}

Management needs visibility into total mess consumption, revenue and billing,
student satisfaction, food wastage, complaint trends, and overall mess
efficiency.

% ============================================================
% 3. CURRENT JOURNEY & LANDSCAPE
% ============================================================

\section{Current Journey \& Landscape}

\subsection*{Student --- current}

\begin{center}
\begin{tabular}{c}
Student joins hostel\\
$\downarrow$\\
Fixed mess fee charged\\
$\downarrow$\\
Student consumes a variable number of meals\\
$\downarrow$\\
No direct link between usage and payment\\
$\downarrow$\\
Student still bears the fixed cost
\end{tabular}
\end{center}

\textbf{Core problem}

The payment model does not reflect actual meal consumption.

\subsection*{Mess operations --- current}

\begin{center}
\begin{tabular}{c}
Estimate attending students\\
$\downarrow$\\
Prepare food\\
$\downarrow$\\
Students arrive\\
$\downarrow$\\
Actual demand differs from estimate\\
$\downarrow$\\
Excess or insufficient food\\
$\downarrow$\\
Wastage or shortage
\end{tabular}
\end{center}

The two problems are directly connected: poor consumption visibility leads to
both inaccurate billing and inaccurate food planning.

\subsection*{Existing alternatives}

\begin{center}
\small
\begin{tabularx}{0.98\textwidth}{L{0.23\textwidth} L{0.28\textwidth} Y}
\toprule
\textbf{Alternative} & \textbf{Advantage} & \textbf{Limitation}\\
\midrule
Fixed mess fee & Simple administration & Student pays regardless of usage\\
Manual attendance & Low setup cost & Difficult to maintain accurately\\
Paper registers & Simple record-keeping & Difficult to analyze\\
WhatsApp & Easy communication & No structured consumption data\\
Google Forms & Easy feedback collection & Not connected to billing or operations\\
Excel / Sheets & Flexible & Heavy manual management\\
\bottomrule
\end{tabularx}
\end{center}

% ============================================================
% 4. PROPOSED SOLUTION
% ============================================================

\section{Proposed Solution}

\subsection*{Elevator pitch}

SmartMess is a digital mess-management and consumption platform that connects
what a student eats with how the mess operates. Students view meals, record
consumption, track usage, and give feedback; administrators manage billing,
monitor demand, and reduce food wastage.

The long-term goal is to move from a fixed-fee mess model to a fairer,
consumption-aware model --- students are charged for eligible meals consumed
rather than a blanket amount --- while the same consumption data helps the mess
prepare food closer to actual demand.

\subsection*{Top MVP value propositions}

\begin{enumerate}
    \item \textbf{Fairer meal payments}\\
    Students pay according to eligible meals consumed instead of automatically
    paying for unused meals --- a particular benefit for students who travel,
    skip meals, or are frequently off campus.

    \item \textbf{Food waste reduction}\\
    The mess gets better visibility into expected meal participation and can
    make more informed preparation decisions. Less uncertainty leads to better
    planning and less excess food.

    \item \textbf{Administrative control}\\
    Administrators get one system to manage students, meals, consumption,
    billing, menus, feedback, complaints, and reports.

    \item \textbf{Better student experience}\\
    Students can see what is being served, what they consumed, what they are
    being charged for, what they have paid, and how to give feedback --- creating
    transparency between students and mess management.
\end{enumerate}

\subsection*{Conceptual model}

Student meal selection drives two downstream flows that converge in one admin
system:

\begin{center}
\begin{tabularx}{0.92\textwidth}{L{0.43\textwidth} Y}
\toprule
\textbf{Consumption $\rightarrow$ Billing} & \textbf{Consumption $\rightarrow$ Planning}\\
\midrule
Select meal $\rightarrow$ Meal consumption recorded $\rightarrow$
Student bill $\rightarrow$ Fairer payment &
Select meal $\rightarrow$ Meal consumption recorded $\rightarrow$
Food demand signal $\rightarrow$ Food planning\\
\bottomrule
\end{tabularx}
\end{center}

Both flows feed the admin system, which spans menu management, feedback, and
analytics --- driving a better overall mess experience.

% ============================================================
% 5. GOALS & SUCCESS METRICS
% ============================================================

\section{Goals \& Success Metrics}

The MVP should validate three hypotheses. Targets below are illustrative
starting points --- replace with figures agreed with the pilot institution
before launch.

\begin{center}
\small
\begin{tabularx}{0.99\textwidth}{L{0.27\textwidth} L{0.37\textwidth} Y}
\toprule
\textbf{Hypothesis} & \textbf{How we'll know} & \textbf{Illustrative target}\\
\midrule
\textbf{Payment efficiency:}\\
students can reconcile eligible consumption against charges &
Share of students who view their billing breakdown at least weekly;
support tickets about billing disputes &
$\geq$60\% weekly viewers; $<$5\% of students file a billing dispute per cycle\\
\addlinespace
\textbf{Food efficiency:}\\
participation data improves preparation and cuts waste &
Variance between expected and actual meal participation; reported food wastage &
Forecast error reduced by $\geq$20\% vs. current estimation method\\
\addlinespace
\textbf{Administrative efficiency:}\\
manual processes are replaced by one system &
Admin time spent on billing/consumption reconciliation; adoption of digital
menu \& feedback tools &
$\geq$50\% reduction in manual reconciliation time per billing cycle\\
\bottomrule
\end{tabularx}
\end{center}

% ============================================================
% 6. MVP / FUNCTIONAL REQUIREMENTS
% ============================================================

\section{MVP / Functional Requirements}

Requirements are organized around the most important Critical User Journeys
(CUJs) and prioritized:

\begin{itemize}
    \item \priority{P0} Required for MVP
    \item \priority{P1} Important for a minimum-delightful product
    \item \priority{P2} Future enhancement
\end{itemize}

\subsection{CUJ 1 --- Student manages meal consumption}

\textbf{Goal:} a student can clearly indicate and track the meals they intend
to consume.

\begin{itemize}
    \item \priority{P0} Student can log in.
    \item \priority{P0} Student can view available meals.
    \item \priority{P0} Student can mark meal participation.
    \item \priority{P0} Student can opt out of a meal before the configured cutoff.
    \item \priority{P0} System records eligible meal consumption.
    \item \priority{P1} Student can view meal history.
    \item \priority{P1} Student can see total meal usage for a selected period.
\end{itemize}

\subsection{CUJ 2 --- Student understands and manages payment}

\textbf{Goal:} students have transparency over what they are being charged for.

\begin{itemize}
    \item \priority{P0} Student can view applicable meal rates.
    \item \priority{P0} Student can view eligible meals consumed.
    \item \priority{P0} Student can view the corresponding accumulated amount.
    \item \priority{P0} Student can view current/unpaid mess balance.
    \item \priority{P1} Student can view previous billing periods.
    \item \priority{P1} Student can view a breakdown of charges by meal/date.
    \item \priority{P2} Online payment integration.
    \item \priority{P2} Refund/credit handling for eligible skipped meals.
\end{itemize}

The billing model must be configurable per institution, since colleges differ in
mess-fee policy.

\subsection{CUJ 3 --- Student checks menu \& gives feedback}

\textbf{Goal:} students know what they are eating and can influence future food
quality.

\begin{itemize}
    \item \priority{P0} Student can view the daily menu.
    \item \priority{P0} Student can view upcoming meals.
    \item \priority{P1} Student can rate meals.
    \item \priority{P1} Student can submit written feedback.
    \item \priority{P1} Student can report food-quality issues.
    \item \priority{P2} Student can provide category-specific ratings for taste,
          quality, quantity, and hygiene.
\end{itemize}

\subsection{CUJ 4 --- Administrator manages billing \& consumption}

\textbf{Goal:} administrators manage consumption-based records without separate
manual systems.

\begin{itemize}
    \item \priority{P0} Admin can manage student accounts.
    \item \priority{P0} Admin can configure meal rates.
    \item \priority{P0} Admin can view student meal participation.
    \item \priority{P0} Admin can view recorded consumption.
    \item \priority{P0} Admin can view student billing information.
    \item \priority{P0} Admin can generate billing summaries.
    \item \priority{P1} Admin can adjust eligible records according to
          institutional policies.
    \item \priority{P1} Admin can export billing/consumption reports.
\end{itemize}

\subsection{CUJ 5 --- Administrator manages menu \& feedback}

\textbf{Goal:} administrators control what is served and understand student
response.

\begin{itemize}
    \item \priority{P0} Admin can create and publish menus.
    \item \priority{P0} Admin can edit menus.
    \item \priority{P0} Admin can view student feedback.
    \item \priority{P1} Admin can identify poorly rated meals.
    \item \priority{P1} Admin can identify recurring food complaints.
    \item \priority{P1} Admin can publish menu changes and announcements.
\end{itemize}

\subsection{CUJ 6 --- Administrator reduces food wastage}

\textbf{Goal:} use meal participation data to improve food preparation decisions.

\begin{itemize}
    \item \priority{P0} Admin can view expected meal participation.
    \item \priority{P0} Admin can compare expected and actual participation.
    \item \priority{P1} Admin can view meal participation trends.
    \item \priority{P1} Admin can identify meals with consistently low participation.
    \item \priority{P2} System estimates expected meal demand using historical data.
    \item \priority{P2} System tracks food prepared versus food consumed/wasted.
\end{itemize}

\subsection{CUJ 7 --- Administrative access \& control}

\textbf{Goal:} only authorized personnel can access sensitive student, billing,
and operational information.

\begin{itemize}
    \item \priority{P0} Role-based access for students and administrators.
    \item \priority{P0} Admin can manage authorized users.
    \item \priority{P0} Students cannot access administrative functions.
    \item \priority{P0} Billing and consumption data is accessible only to
          authorized roles.
    \item \priority{P1} Different administrative permissions can be assigned to
          different staff members.
\end{itemize}

\subsection{Non-functional requirements}

Not covered by the CUJs above but required for a production-ready MVP:

\begin{center}
\small
\begin{tabularx}{0.99\textwidth}{L{0.22\textwidth} Y}
\toprule
\textbf{Area} & \textbf{Requirement}\\
\midrule
Performance &
Meal opt-in/opt-out actions confirm within 2 seconds under normal load;
menu and billing views load within 3 seconds.\\
Availability &
Core student-facing features (view menu, mark participation, view balance)
available during meal windows with minimal planned downtime.\\
Data privacy &
Student billing and consumption records are personal data; access is role-gated
and logged; retention policy defined per institution.\\
Auditability &
Billing calculations and admin adjustments to consumption records are logged
with actor, timestamp, and reason.\\
Scalability &
System supports the largest pilot hostel's population (define target, e.g.
2,000+ students) without redesign.\\
Device support &
Student-facing flows work on mobile web at minimum, given meal-time usage
patterns.\\
\bottomrule
\end{tabularx}
\end{center}

% ============================================================
% 7. USER / OBJECT LIFECYCLE
% ============================================================

\section{User / Object Lifecycle}

\begin{center}
\small
\begin{tabularx}{0.98\textwidth}{L{0.20\textwidth} Y}
\toprule
\textbf{Object} & \textbf{Lifecycle}\\
\midrule
Student & Register $\rightarrow$ Consume meals $\rightarrow$ Track usage
$\rightarrow$ View charges $\rightarrow$ Continue / Deactivate\\
Meal & Created $\rightarrow$ Published $\rightarrow$ Student participation
$\rightarrow$ Consumed $\rightarrow$ Billed $\rightarrow$ Archived\\
Billing & Meal consumption $\rightarrow$ Charge calculation $\rightarrow$ Bill
$\rightarrow$ Payment $\rightarrow$ Transaction history\\
Food & Demand estimate $\rightarrow$ Preparation $\rightarrow$ Consumption
$\rightarrow$ Remaining food $\rightarrow$ Waste measurement\\
Feedback & Meal $\rightarrow$ Rating $\rightarrow$ Feedback $\rightarrow$
Analysis $\rightarrow$ Operational improvement\\
\bottomrule
\end{tabularx}
\end{center}

% ============================================================
% 8. MVP BOUNDARIES
% ============================================================

\section{MVP Boundaries}

\subsection*{Included in MVP}

\begin{itemize}
    \item Student authentication
    \item Administrative authentication
    \item Meal/menu management
    \item Meal participation
    \item Meal opt-out
    \item Consumption tracking
    \item Meal-rate management
    \item Student billing visibility
    \item Basic billing calculation
    \item Feedback and ratings
    \item Complaint management
    \item Food-demand visibility
    \item Basic administrative analytics
\end{itemize}

\subsection*{Future (out of MVP scope)}

\begin{itemize}
    \item Online payment gateway
    \item Automated refunds/credits
    \item QR-based meal verification
    \item Advanced demand prediction
    \item Actual kitchen waste tracking
    \item Nutrition tracking
    \item Multi-hostel/multi-mess support
    \item ERP integration
\end{itemize}

% ============================================================
% 9. MVP PHASES & RELEASE CRITERIA
% ============================================================

\section{MVP Phases \& Release Criteria}

\subsection*{Phase 1 --- Consumption \& Payment Transparency}

\textbf{Scope:} student accounts, menu, meal participation, opt-out,
consumption tracking, meal rates, student billing view.

\textbf{Objective:} validate whether students value and use consumption-based
mess transparency.

\textbf{Release criteria:} opt-in/opt-out flows are reliable under real
meal-window load; billing view matches manual reconciliation for a full cycle;
no critical data-integrity defects.

\subsection*{Phase 2 --- Administration \& Feedback}

\textbf{Scope:} admin dashboard, billing management, student management,
feedback, ratings, complaints, menu management.

\textbf{Objective:} validate whether the system provides meaningful
administrative value.

\textbf{Release criteria:} admins can run a full billing cycle without falling
back to spreadsheets; feedback and complaint workflows are in active use by
mess staff.

\subsection*{Phase 3 --- Food Efficiency}

\textbf{Scope:} demand analytics, participation trends, expected vs. actual
meals, waste indicators, operational reports.

\textbf{Objective:} validate whether consumption data can reduce avoidable food
preparation and wastage.

\textbf{Release criteria:} forecast-to-actual variance is measurably lower than
the pre-SmartMess baseline over at least one full term.

% ============================================================
% 10. ASSUMPTIONS & RISKS
% ============================================================

\section{Assumptions \& Risks}

\begin{center}
\small
\begin{tabularx}{0.99\textwidth}{L{0.36\textwidth} Y}
\toprule
\textbf{Assumption / Risk} & \textbf{Mitigation}\\
\midrule
Institution is willing and able to move to a consumption-based (vs. flat)
billing model &
Keep billing model fully configurable per institution; support flat-fee mode
as a fallback\\
Students reliably mark participation before the cutoff, or the signal is
unreliable &
Set sensible default cutoffs; consider low-friction defaults (e.g. opt-out
rather than opt-in) per institution preference\\
Mess staff adopt digital menu/consumption workflows over paper/WhatsApp habits &
Phase 2 focuses on replacing existing manual admin workflows one at a time,
not all at once\\
Connectivity/device access in hostel dining areas may be inconsistent &
Design for degraded connectivity; confirm minimum device requirement with
pilot site\\
Billing disputes if students forget to opt out and are still charged &
Transparent, itemized billing view (CUJ 2) and clear cutoff communication\\
\bottomrule
\end{tabularx}
\end{center}

% ============================================================
% 11. TELEMETRY / MEASUREMENT
% ============================================================

\section{Telemetry / Measurement}

The system should measure:

\begin{itemize}
    \item Meals offered
    \item Meals opted into / opted out
    \item Meals consumed
    \item Student billing amount
    \item Average cost per meal
    \item Total mess consumption
    \item Feedback submissions
    \item Average meal rating
    \item Complaint frequency
    \item Expected vs. actual participation
    \item Food wastage indicators
\end{itemize}

\subsection*{Key questions}

\begin{itemize}
    \item Are students paying more fairly?
    \item Is the mess preparing food more efficiently?
    \item Is administration gaining better visibility?
\end{itemize}

% ============================================================
% 12. APPENDIX
% ============================================================

\section{Appendix}

Detailed supporting documents referenced by, but not contained in, this PRD:

\subsection*{Product / UX}

\begin{itemize}
    \item Student journey
    \item Admin journey
    \item Payment journey
    \item Wireframes
    \item UI designs
\end{itemize}

\subsection*{Technical}

\begin{itemize}
    \item Database schema
    \item Billing logic
    \item API specification
    \item Authentication architecture
    \item Payment integration
\end{itemize}

\subsection*{Validation}

\begin{itemize}
    \item Student interviews
    \item Mess administrator interviews
    \item Pilot results
    \item Payment-model validation
    \item Food-wastage baseline
\end{itemize}

\subsection*{Business}

\begin{itemize}
    \item Pricing/billing model
    \item College onboarding
    \item Revenue model
    \item Competitive analysis
\end{itemize}

\subsection*{Product Summary}

\textbf{Problem} Student can pay a fixed mess fee despite consuming fewer
meals, while mess operators may prepare food without accurate visibility into
actual demand.

\textbf{Solution} SmartMess connects meal consumption, billing, administration,
feedback, and food-efficiency data into one platform.

\textbf{Primary value} Students: payment fairness and transparency. Mess:
demand visibility and reduced wastage. Administration: centralized control over
consumption, billing, menu, feedback, and operations.

\textbf{Vision} Move hostel mess management from fixed fees + estimated demand +
manual management, to consumption-based transparency + data-driven planning +
centralized administration.

\textbf{SmartMess in one line}

SmartMess makes mess payments fairer for students and mess operations smarter,
by connecting what students consume with what the mess prepares and charges.

\end{document}
